\section*{第 \ref{ch:b1:erosion-transport} 章——侵蚀与搬运}

\begin{solution}{exo:b1:erosion-transport:1}
\qty{0.4}{g.L^{-1}} 即 \qty{0.4}{kg.m^{-3}}，故通量为 $500 \times 0.4 =
\qty{200}{kg.s^{-1}}$。一年以 \num{365.25}~天计，合
\qty{3.156e7}{s}，故年通量为 \qty{6.31e9}{kg}，即
\qty{6.31}{Mt}。
\end{solution}

\begin{solution}{exo:b1:erosion-transport:2}
$w_s = R g d^2/(18\nu) = 1.6548 \times 9.81 \times (\num{30e-6})^2/(18 \times
\num{1.0034e-6}) = \qty{8.09e-4}{m.s^{-1}}$，每秒约 \qty{0.8}{mm}。其
颗粒雷诺数为 $w_s d/\nu = \num{8.09e-4} \times
\num{30e-6}/\num{1.0034e-6} = 0.024$，远小于 1，适用斯托克斯定律。它
下沉 \qty{2}{m} 所需时间为 $2/\num{8.09e-4} = \qty{2472}{s}$，约
\qty{41}{min}。
\end{solution}

\begin{solution}{exo:b1:erosion-transport:3}
$\tau_b = \rho g h S = 998.2 \times 9.81 \times 0.5 \times \num{4e-4} =
\qty{1.96}{Pa}$。取 $(\rho_s - \rho) g = 1651.8 \times 9.81 =
\qty{16204}{N.m^{-3}}$。对 \qty{1}{mm}，$\theta = 1.96/(16204 \times 0.001) =
0.121$，约为\omterm{def:b1:erosion-transport:threshold}{起动阈值} \num{0.0314} 的四倍，故沙粒运动；对
\qty{10}{mm}，$\theta = 0.0121$，为\omterm{def:b1:erosion-transport:threshold}{起动阈值} \num{0.0556} 的五分之一，故
砾石留在原处。水流携带其中的沙，越过一层不动的砾石层。
\end{solution}

\begin{solution}{exo:b1:erosion-transport:4}
(a)~$12.6/3.9 = 3.2$：河流输运的\omterm{def:b1:erosion-transport:load}{悬移质}
约为\omterm{def:b1:erosion-transport:load}{溶解质}的三倍。(b)~大型储库：$0.20 \times 14 = \qty{2.8}{Gt}$ 每
年；全部储库：$0.26 \times 14 = \qty{3.64}{Gt}$ 每年。(c)~溶解
通量每年为 \qty{3.9e15}{g}，浓度为 \qty{120}{g.m^{-3}}，故携带它的水
体积为 $\num{3.9e15}/120 = \qty{3.25e13}{m^3}$，约 \qty{32500}{km^3} 每年。
该水中的\omterm{def:b1:erosion-transport:load}{悬移质}含量为 $\qty{1.26e16}{g}/\qty{3.25e16}{L} =
\qty{0.39}{g.L^{-1}}$。
\end{solution}

\begin{solution}{exo:b1:erosion-transport:5}
(a)~断面面积 $10 \times 2 = \qty{20}{m^2}$，湿周
为 $10 + 2 \times 2 = \qty{14}{m}$，故 $R_h = 20/14 = \qty{1.43}{m}$。
(b)~$\tau_b = 998.2 \times 9.81 \times 1.43 \times 0.001 = \qty{14.0}{Pa}$，
与很宽、深 \qty{2}{m} 的河道的 $\qty{19.6}{Pa}$ 相比，
比值为 $R_h/h = 0.714$，其余摩擦由河岸承担。(c)~对宽为 $b$ 的
矩形河道，$R_h = bh/(b + 2h) = h/(1 + 2h/b)$，
其趋于 $h$，这发生在 $b$ 增大时。
\end{solution}

\begin{solution}{exo:b1:erosion-transport:6}
(a)~$d_1 = \left(18\nu^2/(Rg)\right)^{1/3} = (18/(Rg))^{1/3}\,\nu^{2/3}$，且 $R$
$g$ 均与黏度无关。(b)~$\nu = \num{1.3059e-3}/999.7 =
\qty{1.306e-6}{m^2.s^{-1}}$，\qty{10}{\celsius} 时，与
\qty{20}{\celsius} 时的 \qty{1.0034e-6}{m^2.s^{-1}} 相比，比值为 \num{1.302}，
则 $1.302^{2/3} = 1.192$，故 $d_1 = 0.104 \times 1.192 = \qty{0.124}{mm}$。
(c)~在斯托克斯区间内，$w_s$ 与 $\nu$ 成反比：颗粒
沉降慢 $1.302$ 倍，约 \qty{30}{\percent}。冷水使细颗粒
更久地悬浮：冬季河流携带的粉砂
在相同流量下可被带得稍远一些。
\end{solution}

\begin{solution}{exo:b1:erosion-transport:7}
取 $R g/(18\nu) = \qty{8.99e5}{m^{-1}.s^{-1}}$，并按综合公式：
\begin{center}
\begin{tabular}{@{}lccc@{}}
\hline
$d$ & 斯托克斯 (\si{m.s^{-1}}) & 综合公式 (\si{m.s^{-1}}) & 比值 \\
\hline
\qty{0.1}{mm} & \num{8.99e-3} & \num{7.53e-3} & \num{1.19}\\
\qty{0.5}{mm} & \num{0.225} & \num{0.0711} & \num{3.16}\\
\qty{2}{mm} & \num{3.60} & \num{0.197} & \num{18.3}\\
\hline
\end{tabular}
\end{center}
比值随粒径增大而增大，因为颗粒雷诺数随之增大：斯托克斯
定律仅适用于小雷诺数，按 $d^2$ 增长，而尾流阻力
使粗颗粒速度仅随 $\sqrt d$ 增长。
\end{solution}

\begin{solution}{exo:b1:erosion-transport:8}
(a)~$(Rg/\nu^2)^{1/3} = \qty{25263}{m^{-1}}$。对 \qty{0.2}{mm}，$D_* = 5.05$
且 $\theta_c = 0.3/7.06 + 0.055(1 - e^{-0.101}) = 0.0425 + 0.0053 =
0.0478$；$\tau_c = 0.0478 \times 16204 \times \num{2e-4} = \qty{0.155}{Pa}$。
对 \qty{20}{mm}，$D_* = 505$，且 $\theta_c = 0.3/607 + 0.055(1 - e^{-10.1}) =
0.0005 + 0.0550 = 0.0555$；$\tau_c = 0.0555 \times 16204 \times 0.02 =
\qty{18.0}{Pa}$。(b)~$h = \tau_c/(\rho g S)$，其中 $\rho g S =
\qty{19.58}{Pa.m^{-1}}$：\qty{7.9}{mm} 水深即可使细沙运动，\qty{0.92}{m}
水深方能使砾石运动。
\end{solution}

\begin{solution}{exo:b1:erosion-transport:9}
流域每年损失 \qty{0.8}{Mt} = \qty{8e5}{t}，即 $\num{8e5}/2000 =
\qty{400}{t.km^{-2}.\text{年}^{-1}}$。由\cref{prop:b1:erosion-transport:denudation}可得
$D = \qty{8e8}{kg}/(2650 \times \qty{2e9}{m^2}) = \qty{1.51e-4}{m}$ 每年，或
\qty{151}{m/Ma}。这约为以铍-10 测得流域的中位数
\qty{54}{m/Ma} 的三倍，但低于其平均值 \qty{218}{m/Ma}：
这是一个相当快的流域，符合陡峭或构造活跃的
地区的特点。
\end{solution}

\begin{solution}{exo:b1:erosion-transport:10}
(a)~$u_* = \sqrt{g h S} = \sqrt{9.81 \times 3 \times \num{5e-4}} =
\qty{0.1213}{m.s^{-1}}$，且 $\tau_b = \rho u_*^2 = \qty{14.7}{Pa}$。
(b)~悬浮要求 $w_s < u_* = \qty{0.1213}{m.s^{-1}}$。综合公式
给出 \qty{1}{mm} 处的 \qty{0.126}{m.s^{-1}}；用试算法
解 $w_s(d) = 0.1213$，得 $d = \qty{0.94}{mm}$。(c)~对砾石
$\theta_c \approx 0.0555$，故 $\tau_c = \qty{14.7}{Pa}$ 对应 $d \approx
14.7/(0.0555 \times 16204) = \qty{16.3}{mm}$；该粒径下阈值曲线给出
$\theta_c = 0.0555$，确与之相符。(d)~粒径大于约 \qty{16}{mm} 的颗粒
留在床面上，形成稳定的砾石层；较细的砾石及粗于
约一毫米的沙以\omterm{def:b1:erosion-transport:load}{推移质}运动；更细的沙、粉砂与黏土
则以\omterm{def:b1:erosion-transport:load}{悬移质}运动。
\end{solution}

\begin{solution}{exo:b1:erosion-transport:11}
(a)~$0.30 \times \qty{1000}{m} = \qty{300}{m}$ 的岩石，对整个
区域取平均。(b)~回弹量为 $300 \times 2700/3300 = \qty{245}{m}$。(c)~平均
地表仅降低 $300 - 245 = \qty{55}{m}$。(d)~未被侵蚀的山顶，
随整个区域一同抬升：它们现在比原来
高出 \qty{245}{m}。切割河谷可以抬高山脉的峰顶。
\end{solution}

\begin{solution}{exo:b1:erosion-transport:12}
(a)~$0.9 \times \qty{5e8}{kg}/\qty{1300}{kg.m^{-3}} = \qty{346000}{m^3}$ 每年，
即占库容的 \qty{0.69}{\percent}。(b)~半库容，
\qty{25e6}{m^3}，按此速率充满需 $\num{25e6}/\num{346000} = \num{72}\text{ 年}$。
(c)~随水库淤积，其库容及水的停留时间
随之减小，较细的颗粒得以穿过：其捕获效率下降。
同时，较老的沉积物被埋藏压实，挤出水分，使
其体积密度升高，每吨沉积物占据的空间变小。两种效应共同
减缓了水库老化过程中的库容损失。
\end{solution}

\begin{solution}{pb:b1:erosion-transport:1}
\textbf{1.} 一年的水量为 $1.2 \times \num{3.156e7} = \qty{3.79e7}{m^3}$；
均摊到 \qty{8e7}{m^2} 的面积上，水层厚 \qty{0.473}{m}：每年的径流深为
\qty{473}{mm}。

\textbf{2.} $1.2 \times \qty{0.150}{kg.m^{-3}} = \qty{0.18}{kg.s^{-1}}$，即
每年 \qty{5680}{t}。

\textbf{3.} $1.2 \times 0.080 = \qty{0.096}{kg.s^{-1}}$，即
每年 \qty{3030}{t}。

\textbf{4.} \omterm{def:b1:erosion-transport:load}{推移质}为每年 \qty{568}{t}，总量为 $5680 + 3030 +
568 = \qty{9278}{t}$ 每年。

\textbf{5.} \omterm{def:b1:erosion-transport:load}{悬移质}浓度在洪水的少数几天很高，
此时流量也大，其余时间则两者皆低。时间平均
对每一天赋予相同权重，因此结果主要由
流量小、浓度低的众多日子决定。由
\cref{prop:b1:erosion-transport:flux}可知，年输沙量为
$\int Q\,C\,\mathrm dt$；流量加权平均即为该积分
除以 $\int Q\,\mathrm dt$，故应将流量加权平均值乘以
平均流量。

\textbf{6.} $3030/9278 = \num{0.33}$，即三分之一。对全球河流而言，
$3.9/(12.6 + 3.9) = \num{0.24}$，若计入\omterm{def:b1:erosion-transport:load}{推移质}，则比值更小。
本流域的比例更高：丘陵地区机械侵蚀缓慢，
岩石有充足时间溶解，石灰岩与白垩
则完全溶解。

\textbf{7.} $9278/80 = \qty{116}{t.km^{-2}.\text{年}^{-1}}$。

\textbf{8.} $D = \qty{9.278e6}{kg}/(2700 \times \qty{8e7}{m^2}) =
\qty{4.30e-5}{m}$ 每年，即 \qty{43.0}{m/Ma}。

\textbf{9.} $43.0/54 = 0.80$：略低于实测流域的中位数，
也比其平均值低五倍。这是一个中等偏慢的流域，符合
平静地区丘陵的预期。

\textbf{10.} 实测径流量只平均了少数几年的测量期，且取决于其中
发生的洪水。铍-10 平均的是侵蚀掉产生它的几米岩层所需的时间：
在 \qty{43}{m/Ma} 下，每米需
\num{23000} 年，因此平均时段为数万年。二者可能不同，原因包括
短期记录中缺失的罕见大洪水，
近几个世纪的农耕与开垦加剧了侵蚀，
以及沉积物在河漫滩中被储存并再次带走。

\textbf{11.} \omterm{def:b1:erosion-transport:load}{溶解质}的一部分并非来自岩石：其中大部分
碳酸氢根携带的碳取自空气或土壤，雨水还带来海
盐。若将其计入，就等于把流域并未损失的岩石也算了进去。

\textbf{12.} $\num{1.5e6}/1.2 = \qty{1.25e6}{s}$，即 \num{14.5} 天。

\textbf{13.} 综合公式给出 \qty{7.53e-3}{m.s^{-1}}：颗粒下落
\qty{8}{m} 需 \qty{1062}{s}，约 \qty{18}{min}。颗粒被捕获。

\textbf{14.} $w_s = \num{8.99e5} \times (\num{2e-5})^2 =
\qty{3.60e-4}{m.s^{-1}}$：\qty{8}{m} 需 \qty{22200}{s}，约 \num{6.2}~小时。
颗粒被捕获。

\textbf{15.} $w_s = \num{8.99e5} \times (\num{1e-6})^2 =
\qty{8.99e-7}{m.s^{-1}}$：\qty{8}{m} 需 \qty{8.9e6}{s}，即 \num{103}~天，
是停留时间的七倍。大部分颗粒随水流穿过大坝。

\textbf{16.} 颗粒须在 \qty{1.25e6}{s} 内下落 \qty{8}{m}，
速度至少为 $\qty{6.4e-6}{m.s^{-1}}$。由斯托克斯定律 $d = \sqrt{18\nu
w_s/(Rg)} = \sqrt{\num{6.4e-6}/\num{8.99e5}} = \qty{2.7}{\micro m}$，属细
黏土。较粗的颗粒在任何进入之处均被捕获。

\textbf{17.} $568 + 0.7 \times 5680 = \qty{4544}{t}$ 每年，
可充填 $\num{4.544e6}/1300 = \qty{3496}{m^3}$，占水库
库容的 \qty{0.233}{\percent}。一半即 \qty{7.5e5}{m^3}，约 \num{215} 年充满。

\textbf{18.} $100/43.0 = \qty{2.33}{Ma}$。

\textbf{19.} $1 - 2700/3300 = 0.18$：地表仅降低
所去除厚度的 \qty{18}{\percent}。

\textbf{20.} $100/0.1818 = \qty{550}{m}$ 的岩石。

\textbf{21.} 整个区域因回弹抬升 $550 \times 2700/3300 = \qty{450}{m}$，
几乎未被侵蚀的山峰随之抬升，最高可达
\qty{450}{m}，同时山谷加深：地势起伏先增大，而后减小。

\textbf{22.} \omterm{def:b1:erosion-transport:denudation}{剥蚀速率}首先取决于坡度。随着丘陵
被降低、坡度变缓，速率将减小，除非隆升或
气候变化使其重新变陡：按当前速率算得的时间为
下界。

\textbf{23.} $550/42.95 = 12.80$。按当前速率并计入均衡回弹，
将瓦尔登丘陵降低 \qty{100}{m} 需 \textbf{12.8}~百万年。
\end{solution}
